\section{Results}
\label{sec:results}

All numbers are from the 77-clip LAFAN1 corpus (34 floor-contact / 43
locomotion clips) evaluated with the reference-free metric suite
(Section~\ref{sec:grounding}); the multi-stage pipeline
(Section~\ref{sec:results-multistage}) and \texttt{gmr\_heightfix}
(Section~\ref{sec:results-heightfix}) are evaluated on the complete corpus,
the global pipeline (Section~\ref{sec:results-global}) on the 36/77 clips
completed at time of writing. All comparisons are against
\texttt{gmr\_heightfix} --- GMR's own published constant-offset mitigation
--- not unmodified GMR output, which is strictly worse and reported only as
motivation (Section~\ref{sec:results-motivation}).

\subsection{Motivation}
\label{sec:results-motivation}

Unmodified GMR output on floor-contact clips shows 12.9--15.9~cm maximum
floor penetration affecting 39--91\% of frames, against 1.0~cm and 0.3\% of
frames on GMR's own clean-locomotion clips (\texttt{walk1\_subject1}). This
gap holds at full-corpus scale, not only on individual clips.

\subsection{GMR's Own Mitigation is Insufficient}
\label{sec:results-heightfix}

\texttt{gmr\_heightfix} reduces penetration but leaves the robot floating
well above the floor at exactly the frames that should show tight ground
contact, since one rigid per-clip offset cannot satisfy a standing phase
and a lying phase in the same clip simultaneously
(Table~\ref{tab:heightfix}, \texttt{joint\_ok\%}).

\begin{table}[t]
\centering
\caption{\texttt{gmr\_heightfix}, class means.}
\label{tab:heightfix}
\begin{tabular}{lcc}
\toprule
Metric & Floor ($n=34$) & Locomotion ($n=43$) \\
\midrule
Floor penetration, max (cm)     & 2.76  & 3.06  \\
Floor penetration (\% frames)   & 0.38  & 3.42  \\
Self-collision (\% frames)      & 6.34  & 3.85  \\
Worst float (cm)                & 18.04 & 6.61  \\
Velocity spikes (mean/clip)     & 0.18  & 0.00  \\
Peak joint velocity (rad/s)     & 34.04 & 32.82 \\
Joint contact-correctness (\%)  & 0.19  & 46.26 \\
\bottomrule
\end{tabular}
\end{table}

\subsection{Multi-Stage Pipeline}
\label{sec:results-multistage}

\begin{table}[t]
\centering
\caption{Multi-stage pipeline, class means.}
\label{tab:multistage}
\begin{tabular}{lcc}
\toprule
Metric & Floor ($n=34$) & Locomotion ($n=43$) \\
\midrule
Floor penetration, max (cm)     & 0.00  & 0.00  \\
Floor penetration (\% frames)   & 0.00  & 0.00  \\
Self-collision (\% frames)      & 0.002 & 0.001 \\
Worst float (cm)                & 7.55  & 5.57  \\
Velocity spikes (mean/clip)     & 0.00  & 0.00  \\
Peak joint velocity (rad/s)     & 37.39 & 37.92 \\
Joint contact-correctness (\%)  & 98.85 & 98.79 \\
\bottomrule
\end{tabular}
\end{table}

Five of six metrics improve on \texttt{gmr\_heightfix} by more than an
order of magnitude; joint contact-correctness rises from 0.2\% to 98.8\% on
the floor class. Peak joint velocity is 9.8\% (floor) and 15.5\%
(locomotion) above \texttt{gmr\_heightfix} --- a real cost, attributable to
branch-flip artifacts in the per-frame correction stage
(Section~\ref{sec:discussion}) and only partially mitigated by a null-space
continuity term in that stage; the mitigation itself trades a smaller
regression in jerk and contact-slip on part of the floor class.

\subsection{Oracle Comparison}
\label{sec:results-oracle}

A per-clip constant Z-shift, swept by grid search to jointly minimize
penetration and floating with no notion of per-frame contact timing,
outperforms the multi-stage pipeline on either single-axis metric in
isolation (held-frame float within tolerance, or whole-body penetration,
each considered alone) on some clips. It cannot satisfy joint
contact-correctness, which requires both at the same frame: a rigid shift
cannot correct a multi-phase clip's standing stance and its lying phase
with one number. This is the basis for reporting joint contact-correctness
as the primary contact metric.

\subsection{Global Pipeline}
\label{sec:results-global}

Full-corpus evaluation of the global pipeline (Section~\ref{sec:global}) is
in progress; we report results on the 36 clips completed at time of writing
(18 floor-contact, 18 locomotion, class-balanced), pending completion of the
remaining 41.

\begin{table*}[t]
\centering
\caption{Global pipeline vs.\ \texttt{gmr\_heightfix} and the multi-stage
pipeline, class means, 36/77 clips.}
\label{tab:global}
\begin{tabular}{llccc}
\toprule
Metric & Class & \texttt{gmr\_heightfix} & Multi-stage & Global \\
\midrule
\multirow{2}{*}{Joint contact-correctness (\%)}
  & Floor      & 0.20  & 98.50 & 81.02 \\
  & Locomotion & 46.51 & 99.84 & 81.27 \\
\multirow{2}{*}{Floor penetration (cm)}
  & Floor      & 3.06  & 0.00  & 0.00  \\
  & Locomotion & 3.07  & 0.00  & 0.00  \\
\multirow{2}{*}{Self-collision (\% frames)}
  & Floor      & 6.14  & 0.01  & 1.13  \\
  & Locomotion & 2.87  & 0.00  & 0.91  \\
\multirow{2}{*}{Worst float (cm)}
  & Floor      & 18.35 & 8.16  & 13.54 \\
  & Locomotion & 7.87  & 5.03  & 11.53 \\
\multirow{2}{*}{Joint jerk}
  & Floor      & 5293  & 2601  & 803   \\
  & Locomotion & 7409  & 3101  & 1256  \\
\multirow{2}{*}{Peak joint velocity (rad/s)}
  & Floor      & 33.04 & 31.31 & 13.01 \\
  & Locomotion & 40.11 & 37.32 & 11.67 \\
\multirow{2}{*}{Velocity spikes (mean/clip)}
  & Floor      & 0.06  & 0.00  & 0.00  \\
  & Locomotion & 0.00  & 0.00  & 0.00  \\
\bottomrule
\end{tabular}
\end{table*}

The global pipeline improves on \texttt{gmr\_heightfix} on joint
contact-correctness, floor penetration, self-collision, and velocity spikes
on both classes, and reduces jerk and peak joint velocity relative to
\textit{both} baselines by a wide margin (jerk 3.2--6.6x lower, peak
velocity 2.4--3.4x lower than the multi-stage pipeline). It trails the
multi-stage pipeline on joint contact-correctness (81\% vs.\ 98--99\%) and
self-collision, consistent with a smaller total per-frame solve budget
(Section~\ref{sec:comparing}). Worst float is a mixed result: better than
\texttt{gmr\_heightfix} on the floor class (18.35 $\to$ 13.54~cm) but worse
on locomotion (7.87 $\to$ 11.53~cm), the one metric where the global
pipeline does not uniformly improve on the naive baseline. Which pipeline
is preferable depends on which axis matters more for a given downstream
use; Section~\ref{sec:disc-tradeoff} and the planned policy-training
evaluation (Section~\ref{sec:status-next-steps}) are intended to make that
trade-off decidable rather than qualitative.
